\section{Measuring GKP code stabilizers}

A second application of our CV fault-equivalence is in coming up with good stabilizer measurement circuits (often called \textit{syndrome extraction circuits}) for general GKP codes.
Suppose we have an $m$-qumode GKP code defined by the stabilizer group:
\begin{equation}
    \SSS = \langle \hat{D}(\tbinom{\bsym{q}_1}{\bsym{p}_1}), \ldots, \hat{D}(\tbinom{\bsym{q}_{2m}}{\bsym{p}_{2m}}) \rangle
\end{equation}
Recall that GKP codes are implemented by measuring $\hat{H}_j \bmod 1$ for $1 \leq j \leq 2m$, where $\hat{H}_j = \tbinom{\bsym{q}_j}{\bsym{p}_j} \odot \tbinom{\bsym{\hat{q}}}{\bsym{\hat{p}}}$, hence $\hat{D}(\tbinom{\bsym{q}_j}{\bsym{p}_j}) = e^{-i2\pi\hat{H}}$.
The projector associated with the outcome $a \in \mathbb{R}$ of this measurement is:
\begin{equation}
    \Pi_j \coloneqq \sum_{k \in \ZZ} e^{-i2\pi a k}\hat{D}(\tbinom{\bsym{q_j}}{\bsym{p_j}})^{a k}
\end{equation}
For clarity, let's drop the subscript $j$ temporarily, and call this $\Pi$.
We will need polar coordinates: for $\tbinom{\bsym{q}}{\bsym{p}} = (q_1, \ldots, q_m, p_1, \ldots, p_m)$, define $r_j e^{i\theta_j} = q_j + ip_j$.
We showed in TODO that we can represent this projector in the ZX-calculus as:
\begin{equation}
    \tikzfig{qumodes/focked_up/gkp_measurement/measurement_general_displacement_outcome_a}
\end{equation}
Since we only care about the error correcting properties of this diagram, we can choose the representative of this diagram in the $\zxmodchi$-calculus that has $a = 0$:
\begin{equation}\label{eq:measurement_general_displacement}
    \tikzfig{qumodes/focked_up/gkp_measurement/measurement_general_displacement}
\end{equation}
This diagram features an $m$-legged $0_\square$-labelled spider; this represents a GHZ state in the multimode square GKP code.
%This diagram features an $m$-legged $0_\square$-labelled spider; this represents a \textit{cat state} in the multimode square GKP code.
%A cat state is just a $\ket{\text{GHZ}_m}$ state but in the other basis:
%\begin{equation}
%    \ket{\text{cat}_m} = H^{\otimes m} \ket{\text{GHZ}_m}
%\end{equation}
Preparing this state with a low error rate is tricky.
Supposing we wanted to actually implement this diagram, it would be more realistic to assume that at best we can start with single-qumode square GKP states $\ket{+_\square}$.
So a natural question is: can we find a circuit that's fault-equivalent to the diagram in~\eqref{eq:measurement_general_displacement} above, but which only uses single-qumode GKP states?
In other words, can we find a circuit that measures this GKP stabilizer generator and has the same behaviour under noise as the diagram above (at least for the ZX noise model)?
The following proposition is all we need.

\begin{proposition}\label{prop:GHZ_prep_even_fault_equiv}
    For all real $r \neq 0$ and integer $m \geq 1$, the following diagrams are fault-equivalent:
    \begin{equation}
        \tikzfig{qumodes/focked_up/gkp_measurement/cat_state_comb_0_braces}
        ~~=~~
        \tikzfig{qumodes/focked_up/gkp_measurement/cat_state_comb_0_inductive_braces}
    \end{equation}
    \begin{proof}
        Call these diagrams $D_1 = D_2$.
        The direction $D_2 \mapsto D_1$ is certainly fault-accounting, so it remains to show $D_1 \mapsto D_2$ is too.
        Errors on output errors only are easily accounted for, as are green $\hat{Z}(p)$ errors on internal edges, so just consider an error $F_2$ consisting of red $\hat{X}(q)$ errors on internal edges:
        \begin{equation}
            \tikzfig{qumodes/focked_up/gkp_measurement/cat_state_comb_0_inductive_j_error}
        \end{equation}
        Suppose that this error is undetectable (otherwise it doesn't need accounting for).
        Note that $D_2$ has detecting webs of the following form, for $1 \leq j \leq m$ and all $a_j \in \frac{1}{r}\mathbb{Z}$:
        \begin{equation}
            \tikzfig{qumodes/focked_up/gkp_measurement/cat_state_comb_0_inductive_j_web_detector}
        \end{equation}
        For error $F_2$ above to not be detected by this web, \Cref{prop:detecting_web_detects_error} tells us that we must have $a_j (t_j + u_j - v_j) \in \mathbb{Z}$ for all $a_j \in \frac{1}{r}\mathbb{Z}$:
        \begin{equation}
            \tikzfig{qumodes/focked_up/gkp_measurement/cat_state_comb_0_inductive_j_web_detector_error}
        \end{equation}
        This in fact occurs iff it's true for $a_j = \frac{1}{r}$, since all other values of $a_j$ are just integer multiples of this.
        So we get $t_j + u_j - v_j = r k_j$, for some $k_j \in \ZZ$.
        Using another web, we can in fact show all these $k_j$ are equal.
        Consider the following web, which exists for all $a \in \RR$ and all $2 \leq j \leq m$, with $F_2$ also drawn:
        \begin{equation}
            \tikzfig{qumodes/focked_up/gkp_measurement/cat_state_comb_0_inductive_j_web_pairwise_detector_error}
        \end{equation}
        \Cref{prop:detecting_web_detects_error} tells us $F_2$ is undetected by this web iff:
        \begin{equation}
            \begin{aligned}
                \phantom{\iff}&& a (t_1 + u_1 - v_1) - a(t_j + u_j - v_j) &\in \mathbb{Z} \\
                \iff&& ark_1 - ark_j &\in \ZZ \\
                \iff&& ar(k_1 - k_j) &\in \ZZ \\
            \end{aligned}
        \end{equation}
        And since this holds for all real $a$, this forces $k_1 - k_j = 0$, i.e.\ $k_1 = k_j$.
        So write $k = k_1 = \ldots = k_m$.
        Plug in $t_j = r k_j - u_j + v_j$ and perform a few rewrites:
        \begin{equation}
            \begin{aligned}
                \tikzfig{qumodes/focked_up/gkp_measurement/cat_state_comb_0_inductive_error}
                &~~=~~ \tikzfig{qumodes/focked_up/gkp_measurement/cat_state_comb_0_inductive_error_sub_1} \\[10pt]
                &~~=~~ \tikzfig{qumodes/focked_up/gkp_measurement/cat_state_comb_0_inductive_error_sub_2}
            \end{aligned}
        \end{equation}
        Now note that the $m$ spiders labelled $\chi_{rk}$ are a stabilizer of the comb spider labelled $\indicator{x \eqmod{r} 0}$, so we can remove them, and push the rest out onto the output legs:
        \begin{equation}
                =~~ \tikzfig{qumodes/focked_up/gkp_measurement/cat_state_comb_0_inductive_error_sub_3}
                ~~=~~ \tikzfig{qumodes/focked_up/gkp_measurement/cat_state_comb_0_inductive_error_sub_4}
        \end{equation}
        This error $F_1$ on the output legs only is therefore an equivalent error in $D_1$ of $F_2$.
        So if its weight is less than or equal to that of $F_2$, we're done.
        \begin{align}
            \wt(F_2) &= \sum_{j=1}^m |t_j| + |u_j| + |v_j| \geq \sum_{j=1}^m |t_j| + \sum_{j=1}^m |v_j - u_j| \\
            \wt(F_1) &= \sum_{j=1}^m |v_j| + \sum_{j=1}^m |v_j - u_j|
        \end{align}
        So if $\sum_{j=1}^m |v_j| \leq \sum_{j=1}^m |t_j|$, we're in business.
        Unfortunately we don't have such a guarantee, but all is not lost!
        Had we swapped the roles of the $t_j$ and $v_j$ terms throughout, we would have ended up with a different equivalent error $F_1'$ for $F_2$:
        \begin{equation}
            \tikzfig{qumodes/focked_up/gkp_measurement/cat_state_comb_0_inductive_error}
            ~~=~~
            \tikzfig{qumodes/focked_up/gkp_measurement/cat_state_comb_0_inductive_error_sub_4_alt}
        \end{equation}
        This weights of $F_2$ and $F_1'$ are then:
        \begin{align}
            \wt(F_2) &= \sum_{j=1}^m |t_j| + |u_j| + |v_j| \geq \sum_{j=1}^m |v_j| + \sum_{j=1}^m |t_j + u_j| \\
            \wt(F_1') &= \sum_{j=1}^m |t_j| + \sum_{j=1}^m |t_j + u_j|
        \end{align}
        In this case, we want the other inequality: if $\sum_{j=1}^m |t_j| \leq \sum_{j=1}^m |v_j|$ then $\wt(F_1') \leq \wt(F_2)$.
        Since the inequality has to hold in one direction, and we have an equivalent error of equal or smaller weight in each case, we're done.
    \end{proof}
\end{proposition}

Let's look at this in action.

\begin{example}
    The $E_8$ lattice in $\RR^{8}$ is defined as the span of the columns of the following \textit{generator matrix}:
    \begin{equation}\label{eq:E8_lattice_generator}
        M =
        \begin{pmatrix}
            1 & 0 & 0 & 0 & 0 & 0 & 0 & \frac{1}{2} \\
            -1 & 1 & 0 & 0 & 0 & 0 & 0 & \frac{1}{2} \\
            0 & -1 & 1 & 0 & 0 & 0 & 0 & \frac{1}{2} \\
            0 & 0 & -1 & 1 & 0 & 0 & 0 & \frac{1}{2} \\
            0 & 0 & 0 & -1 & 1 & 0 & 0 & \frac{1}{2} \\
            0 & 0 & 0 & 0 & -1 & 1 & 0 & \frac{1}{2} \\
            0 & 0 & 0 & 0 & 0 & -1 & 1 & \frac{1}{2} \\
            0 & 0 & 0 & 0 & 0 & 0 & -1 & \frac{1}{2}
        \end{pmatrix}
    \end{equation}
    For any integer $d \geq 2$, this lattice gives rise to a four-qumode GKP code encoding four $d$-dimensional qudits.
    Define $M_d = \sqrt{d}M$ and let the lattice $\Gamma_d$ be the span of the columns of $M_d$.
    We'll take $d = 2$, so that this encodes four qubits.
    Take column $j$ of $M_d$, reorder it into $qqpp$ order (\Cref{subsec:phase_space}), and denote it by $\boldsymbol{\xi}_j$.
    The GKP code is defined by the stabilizer group:
    \begin{equation}
        \SSS = \pres{\hat{D}(\boldsymbol{\xi}_1), \ldots, \hat{D}(\boldsymbol{\xi}_n)}
    \end{equation}
    So to implement it, we need to measure a generating set for it -- may as well choose the one we've given above.
    And the generator we're particularly interested in is the one corresponding to the final column of $M_2$:
    \begin{equation}
        \hat{D}(\boldsymbol{\xi}_8) = \hat{D}(\sqrt{2}(\frac{1}{2}, \ldots, \frac{1}{2})^T)
    \end{equation}
    To write the corresponding projector in the ZX-calculus, we need polar coordinates: $\sqrt{2}(\frac{1}{2} + \frac{1}{2}i) = e^{i\pi / 4}$.
    Then the projector is:
    \begin{equation}
        \tikzfig{qumodes/focked_up/gkp_measurement/measurement_e8}
    \end{equation}
    In particular, the given diagram contains a 4-qumode GHZ state.
    We saw in \Cref{prop:GHZ_prep_even_fault_equiv} above how to rewrite this fault-equivalently to a diagram containing only one-legged $0_L$ states:
    \begin{equation}
        \tikzfig{qumodes/focked_up/gkp_measurement/ghz_4}
        ~~=~~
        \tikzfig{qumodes/focked_up/gkp_measurement/ghz_4_2}
        ~~=~~
        \tikzfig{qumodes/focked_up/gkp_measurement/ghz_4_2_1}
    \end{equation}
    TODO continue -- how to post-select away the bad measurement outcomes?
    \begin{equation}
        \tikzfig{qumodes/focked_up/gkp_measurement/ghz_4_2_1_meas1}
    \end{equation}
\end{example}